\BSRAppendix{D}{实施差距、迁移顺序与恢复方案}
\noindent 本附录区分本地已实施与正式目标。009 预约提交映射已在备份后的本地 SQLite 落库；当前共有 37 张物理表。011 已提供本地在线备份、校验和隔离恢复演练命令，不执行生产库替换。以下其他迁移仍是设计方案，每次须在副本验证，再决定是否用于正式环境。

\subsection*{D.1 当前与目标的差异清单}
\begin{BSRTable}
\rowcolor{TableHead}\color{white}\bfseries 领域 & \color{white}\bfseries 当前实现、目标差距及历史保护措施 \\
账号泛化 & 当前 \texttt{users.role} 单表，\texttt{employee\_profiles} 是可选扩展。目标可维持此映射；若拆顾客、技师子表，应先回填同 ID、复核权限与历史订单归属，不重发用户 ID。\\
类别树 & 当前 \texttt{variants.category/subtype} 为文本；目标 \texttt{categories.parent\_id} 自反树。先去重、建立旧文本到类别 ID 的对应，再加 FK 和防环校验；保留旧文本以便回滚。\\
时间化库存 & 当前活跃准备单件有部分唯一索引，租用有复合键；重叠租期主要靠服务端事务检查。目标多实例数据库须设计排他/锁定策略，不能仅靠前端剩余量或普通外键。\\
预约提交映射 & \texttt{reservation\_requests} 已在本地建立，以实际账号、渠道和标识为复合主键；预约 FK 唯一，摘要校验原负载。旧订单无映射且不补造；生产跨实例冲突控制和监控仍须验证。\\
合同明细 & 当前 \texttt{contracts.snapshot\_json/hash} 与预约行保存条款；目标独立 \texttt{contract\_lines}。回填后按旧合同哈希和逐项总额勾稽，旧版本只读，不追溯改签。\\
政策同意 & 当前 \texttt{damage\_consents} 保存损坏文本版本；目标覆盖损坏与超时全文、摘要和主动同意。历史同意不能自动解释为接受新版超时条款，变更需重新勾选。\\
安全验收 & 当前 \texttt{maintenance\_events}、工单 JSON/状态保存验收事实；目标可选 \texttt{safety\_acceptances} 轮次表。保留每次退回、重修和经理决定，不用“最新状态”覆盖历史。\\
费用分项 & 当前逐件超时在 \texttt{rental\_units}，损坏核费在 \texttt{damage\_cases}，结算在 \texttt{damage\_settlements}。目标 LateCharge/DamageCharge 是否独立建表按审计需求决定；旧订单无规则版不得追收。\\
8\% 押金 & 新报价按规格参考价值计算，旧 \texttt{price\_policies.deposit\_nok} 或种子值可能不同。冻结新单行快照与政策版本，历史单显示其原金额；迁移时不得统一覆盖旧单押金。\\
资金与签署 & 当前 \texttt{payments} 是演示回执，合同可记录本人确认/纸质见证；真实第三方支付、退款到账、回调验签、对账及合格电子签署均未接入。接入前须以 provider 类型隔离演示与真实事件。\\
人员提醒 & 当前排班、打卡、提醒已读及审计均有持久表；目标仍需正式岗位任命、资格资料和通知渠道。已读与订单未解决状态分开，改派后按新责任范围重新计算。\\
\end{BSRTable}

\subsection*{D.2 实施顺序与每一步的出口条件}
\begin{BSRTable}
\rowcolor{TableHead}\color{white}\bfseries 阶段 & \color{white}\bfseries 操作、校验与回退门槛 \\
M0 基线冻结 & 取得业务批准的规则版本、数据字典和权限矩阵；备份数据库与附件，记录文件哈希、表/行数、金额汇总、版本分布和可恢复点。没有恢复演练不进入生产迁移。\\
M1 主数据 & 新建类别映射与门店/规格稳定 ID，对旧文本去重并记录无法映射项；在副本回填，查无孤儿 FK、库存号冲突与类别环。异常项人工裁决，不静默删除。\\
M2 合同与同意 & 从不可变合同快照回填明细，校验各版本租金、押金、税费和摘要；新同意表仅承接有证据的版本。旧版合同保留只读；失败恢复到 M0 备份。\\
M3 逐件与维修 & 回填活跃单件绑定、实还时间与当前/归属店；重叠占用、时间缺失、验收轮次不明进入人工核实。添加约束前先消解冲突，不能自动认定单件可租。\\
M4 费用与资金 & 用冻结的旧规则核对费用和退款，不用新 8\% 或新超时政策重算旧单。接入第三方时先在测试环境演练重复回调、超时、对账差异和退款失败恢复。\\
M5 切换验证 & 灰度读写和权限核对；对比订单数、单件状态、合同摘要、未结押金、资金回执与审计链。核心不变量或账务勾稽失败时停止切换，回退到经验证备份。\\
M6 持续运营 & 建立监控、备份周期、恢复演练、故障处理和数据保留制度；按挪威业务日期核对报表。BSR 在安全、性能、合同、价格及退款验收后签字。\\
\end{BSRTable}

\subsection*{D.3 迁移检查点与禁止操作}
\begin{itemize}
\item 恒定主键、库存号、合同号和外部交易引用不得重用；外键迁移后查孤儿、重复和无效状态。
\item 历史价目、合同、规则同意、维修费率及已结算金额只读；新政策从生效时刻起应用，不能为了展示统一而重写旧订单。
\item 单件占用、备货、交付、归还、结算与调拨接收跨表操作必须在事务中一致提交；外部回执通过幂等引用与补偿记录关联。
\item 备份不仅要生成，还要在隔离环境恢复并比对关键总额。本地命令已验证 WAL、BLOB、完整性、外键、散列及关键行数；正式加密、异地保存、周期调度和 RPO/RTO 仍需验收。真实顾客资料、私有受损照片及生产日志不得进入公开代码仓库。
\end{itemize}

\BSRStatus{M0～M6 是目标实施路线。当前文档没有执行任何生产迁移，也没有取得 BSR 政策或技术上线签字。}
